% Tabela 1 | Resultados por nivel, SO com numeros do artigo original (Paper Training - 15 Nov 2025):
% secoes 2.5, 3.2 e 4.4 do original e os graficos originais (Figs. 8 a 10 do original = Figs. 6 a 8 aqui).
% "Pelo menos sete abaixo de 70" e "notas de 3 a 5" foram lidos dos graficos originais; a leitura da
% Fig. 9 confere com a media de 66,90 do texto original (soma 1.405 / 21).
\begin{table}[H]
  \centering
  \caption{Results of the Kirkpatrick-based evaluation by level}
  \label{tab:results-by-level}
  \footnotesize
  \begin{tabularx}{\linewidth}{@{}>{\raggedright\arraybackslash}p{2.6cm} >{\raggedright\arraybackslash}p{4.2cm} >{\raggedright\arraybackslash}X@{}}
    \toprule
    \textbf{Level} & \textbf{Instrument and sample} & \textbf{Result and action} \\
    \midrule
    L1 Reaction & Online questionnaire, seven criteria, 0--10 scale; 10 sessions & Most scores close to 10; content scored about 5 and 4 in two sessions and infrastructure below 8 in one; items below target were followed up \\
    \addlinespace
    L2 Learning & End-of-session test, pass mark 70/100; 21 trainees & Mean 66.90; at least seven trainees below the pass mark; content reviewed with instructors \\
    \addlinespace
    L3 Behavior & Dual questionnaires (trainee and leader), 3 to 12 months after training; 15 pairs & Threshold 70/100; scores below it trigger investigation of relevance or content \\
    \addlinespace
    L4 Results & Customer survey, five categories, 0--5 scale; four airline customers & Ratings between 3 and 5; the ratings of 3 concerned technical documentation and turnaround time \\
    \addlinespace
    Organizational outcomes & Quality management system and quality-cost records; twelve months & 28\% fewer training failures; 35\% fewer regulatory non-conformances; quality costs about USD 1.2 million a year lower \\
    \bottomrule
  \end{tabularx}
  \notafonte{L1 to L4: Kirkpatrick levels. Source: the authors, from the facility's training, quality and customer records. \todo{F2: resultado do Nível 3 (médias das 15 duplas) e contagens que sustentam 28\% e 35\%}}
\end{table}
